% ============================================================
% A-O-02.tex -- Лекция 2 (1.2). Элементарная теория чисел
% ============================================================

\documentclass[serif, aspectratio=169]{beamer}

\usepackage[T2A,T1]{fontenc}
\usepackage[english,russian]{babel}
\usepackage{XCharter}
\usepackage{amsmath}
\usepackage{graphicx}
\usepackage{tikz}
\usetikzlibrary{arrows.meta,positioning,calc,fit,overlay-beamer-styles}
\usepackage{caption}
\usepackage{subcaption}
\usepackage{array}
\usepackage{booktabs}
%\usepackage{fontawesome5}

\newcommand{\sagemark}{%
  \textcolor{green!40!black}{\textbf{\texttt{sage:...}}}%
}

\usetheme{Math}
\usepackage{mymacros}
% \colorscheme{green}
% \colorscheme{terracotta}

% ------------------------------------------------------------
% Title data
% ------------------------------------------------------------
\author{Павел Владимирович Снурницын}
\title{Лекция 2 (1.2)\\ Элементарная теория чисел}
\titlecontext{Инструменты машинно-вспомогательной арифметики\\ с/к, осень 2026}
\institute{к.ф.-м.н., с.н.с., ИБ ВМК}
\date{\small \today}

\setfooterauthor{П.В. Снурницын}
\setfootertopic{1.2 Элементарная теория чисел}
\setfootercontext{Спецкурс A-O 2026}

\newif\ifanimated
\newcommand{\animateon}{\beamerdefaultoverlayspecification{<+-| alert@+>}\animatedtrue}
\newcommand{\animateoff}{\beamerdefaultoverlayspecification{}\animatedfalse}
% TikZ-ключи "step on=<n>" (для текста узлов) и "step draw on=<n>" (для линий/
% стрелок): элемент невидим до шага n, ровно на шаге n подсвечивается цветом
% ThemeAlert (как это делает <+-| alert@+> для списков), после -- чёрный.
% В статичном режиме (\animateoff) -- всегда просто чёрный и видимый.
\newcommand{\steponhelper}[2]{% #1 = свойство цвета (text/draw), #2 = номер шага
  \ifanimated
    \edef\stepon@tmp{alt=<#2>{#1=ThemeAlert}{alt=<#2->{#1=black}{invisible}}}%
  \else
    \edef\stepon@tmp{#1=black}%
  \fi
  \expandafter\pgfkeysalso\expandafter{\stepon@tmp}%
}
\tikzset{
  step on/.code={\steponhelper{text}{#1}},
  step draw on/.code={\steponhelper{draw}{#1}},
}
%\animateon
\animateoff

\begin{document}

% ================================================================
\begin{frame}
    \titlepage
\end{frame}
\note{}

% ============================================================
\section{Делимость}
% ============================================================

% ================================================================
\begin{frame}{Делимость и деление с остатком}
    \begin{itemize}
        \item Определение: $a \mid b$, если $b = aq$ для некоторого целого $q$.
        \item Базовые свойства делимости:
        \begin{itemize}
            \item $a\mid a$ (рефлексивность);
            \item $1\mid a$ и $a\mid 0$ для любого целого $a$;
            \item $a\mid b,\ b\mid c \Rightarrow a\mid c$ (транзитивность);
            \item $a\mid b,\ b\mid a \Rightarrow a =\pm b$ (антисимметричность*);
            \item $a\mid b,\ a\mid c \Rightarrow a \mid (b\pm c)$;
            \item $a\mid b \Rightarrow a\mid bc$ для любого целого $c$;
            \item $a\mid b,\ a\mid c \Rightarrow a \mid (xb+yc)$ для любых целых $x,y$.
        \end{itemize}
        \item Теорема о делении с остатком: для любых $a, b\in\mathbb{Z}$, $b>0$, существуют единственные $q, r$ с $a = bq + r$, $0 \le r < b$.
    \end{itemize}
\end{frame}
\note{%
\begin{itemize}
    \item Свойство $a\mid b,\ b\mid a \Rightarrow a=\pm b$ --- это не полноценная антисимметричность (она давала бы $a=b$), а антисимметричность "с точностью до знака": делимость задаёт частичный порядок на $\mathbb{N}$, а на $\mathbb{Z}$ --- только предпорядок (эквивалентными оказываются $a$ и $-a$)
\end{itemize}
}

% ================================================================
\begin{frame}{НОД, алгоритм Евклида и тождество Безу}
    \begin{itemize}
        \item Наибольший общий делитель (НОД) $d=\gcd(a,b)$ --- натуральное число такое, что
        \begin{itemize}
            \item $d\mid a$ и $d\mid b$;
            \item $c\mid a$ и $c\mid b \Rightarrow c\mid d$.
        \end{itemize}
        \item Теорема (тождество Безу): для любых $a,b$ существуют целые $x,y$, такие что $ax + by = \gcd(a,b)$.
        \item Следствие: если $\gcd(a,b)=1$ (т.е. взаимно просты), то $ax+by=1$ разрешимо.
        \item Множество $\{ax+by : x,y\in\mathbb{Z}\}$ --- в точности все кратные $d = \gcd(a,b)$
        \item Алгоритм Евклида: $\gcd(a,b) = \gcd(b,\, a \bmod b)$ при $b\ne0$, с базой рекурсии $\gcd(a,0)=|a|$.
        %\item \hfill\faTerminal
        \item \sagemark
    \end{itemize}
\end{frame}
\note{%
\begin{itemize}
    \item Алгоритм Евклида --- один из древнейших и самых эффективных алгоритмов в математике (работает за $O(\log \min(a,b))$ шагов)
    \item Идея восстановления коэффициентов Безу --- "раскрутка назад" шагов алгоритма Евклида. Общая схема: $a=bq_0+r_1$, $b=r_1q_1+r_2$, $r_1=r_2q_2+r_3$, $\dots$, $r_{n-2}=r_{n-1}q_{n-1}+r_n$, $r_{n-1}=r_nq_n$ (остаток $0$). Последовательность остатков $b>r_1>r_2>\dots>r_n>0$ строго убывает и ограничена снизу нулём $\Rightarrow$ процедура останавливается за конечное число шагов, и $r_n=\gcd(a,b)$
    \item Пример: $\gcd(93,42)$: $93=2\cdot42+9$, $42=4\cdot9+6$, $9=1\cdot6+3$, $6=2\cdot3+0 \Rightarrow \gcd=3$. Раскручиваем назад: $3=9-1\cdot6=9-1\cdot(42-4\cdot9)=5\cdot9-42=5\cdot(93-2\cdot42)-42=5\cdot93-11\cdot42$, то есть $x=5,\ y=-11$
\end{itemize}
}

% ================================================================
\begin{frame}{Простые числа}
    \begin{itemize}
        \item Число $p>1$ --- простое: делители только $1$ и $p$.
        \item Теорема (Евклид): простых чисел бесконечно много.
        \item Основная теорема арифметики: каждое целое $n>1$ представимо в виде $n = p_1^{a_1}p_2^{a_2}\cdots p_k^{a_k}$, где $p_1<\dots<p_k$ --- простые, $a_i\ge1$, и такое разложение единственно.
        \item Теорема (Адамар, Валле-Пуссен): пусть $\pi(x)$ --- количество простых чисел, не превосходящих $x$, тогда $\pi(x) \sim x/\ln x$ при $x\to\infty$.
        \item Дзета-функция Римана: $\zeta(s) = \displaystyle\sum_{n=1}^{\infty} \frac{1}{n^s}$, $\operatorname{Re}(s)>1$;
        \item Тождество Эйлера: $\zeta(s) = \displaystyle\prod_{p\ \text{простое}} \left(1-p^{-s}\right)^{-1}$.
    \end{itemize}
\end{frame}
\note{%
\begin{itemize}
    \item Идея доказательства бесконечности простых: если бы простых было конечно $p_1,\dots,p_n$, число $N=p_1\cdots p_n+1$ не делится ни на одно из них $\Rightarrow$ у него есть новый простой делитель --- противоречие.
    \item Тождество Эйлера --- аналитическая переформулировка основной теоремы арифметики: раскрывая каждый множитель в геометрическую прогрессию $(1-p^{-s})^{-1}=\sum_{k\ge0}p^{-ks}$ и перемножая такие ряды по всем простым, в силу единственности разложения на простые каждое $n$ встречается в раскрытом произведении ровно один раз --- так и получается сумма $\sum_n n^{-s}$
    \item Именно через $\zeta(s)$ (и, в частности, через то, что она не обращается в ноль на прямой $\operatorname{Re}(s)=1$) доказывается асимптотика $\pi(x)\sim x/\ln x$ из теоремы выше --- это и есть содержательная связь дзета-функции с распределением простых
\end{itemize}
}

% ================================================================
\begin{frame}{Некоторые открытые проблемы о простых}
    \begin{itemize}[<.->]
        \item Гипотеза Гольдбаха: Каждое чётное $n>2$ --- сумма двух простых.
        \item Гипотеза о простых-близнецах: Бесконечно ли пар $(p, p+2)$?
        \item Гипотеза Лежандра: Для любого $n\ge1$ между $n^2$ и $(n+1)^2$ всегда есть простое число.
        \item Простые вида $n^2+1$: Бесконечно ли простых такого вида?
        \item \sagemark
    \end{itemize}
\end{frame}
\note{%
\begin{itemize}
    \item Эти задачи известны как четыре проблемы Ландау. Сформулированы Э. Ландау в 1912 г. на 5-м Международном конгрессе математиков в Кембридже, замечательны тем, что все постановки элементарны, но задачи до сих пор не решены. По двум задачам есть прогресс в ослабленных формулировках.
    \item Тернарная (слабая) проблема Гольдбаха: каждое нечётное $n>5$ --- сумма трёх простых. Гипотеза восходит к переписке Гольдбаха и Эйлера 1742 г.; Харди и Литлвуд в 1923 г. доказали её для достаточно больших $n$ в предположении обобщённой гипотезы Римана; И.М. Виноградов в 1937 г. убрал эту гипотезу, разработав метод оценки тригонометрических сумм по простым числам, и доказал утверждение безусловно --- но лишь для «достаточно больших» $n$ без явной границы. Дальнейшее уточнение эффективных границ (Чен, Лю--Ван и др.) вместе с прямой компьютерной проверкой малых случаев в итоге привело к полному доказательству для всех $n>5$ --- его дал Х.А. Хелфготт в 2013 г., развив метод Харди--Литлвуда--Виноградова
    \item Проблема близнецов не решена, но есть решение в задаче ограниченных разрывов между простыми: в 2013 г. Чжан показал $\liminf(p_{n+1}-p_n) < 7\times10^7$; далее в течение нескольких месяцев бурной коллективной работы (проект Polymath8a) оценка была снижена до $4680$; Мейнард и (независимо) Тао предложили другой, более простой метод, сразу давший $600$; объединив оба подхода, Polymath8b получил текущий безусловный рекорд $246$. При условии дополнительной (обобщённой) гипотезы Эллиотта--Халберстама оценка снижается до $6$; это принципиальный предел этих методов.
    \item Гипотеза Лежандра и проблема простых вида $n^2+1$ остаются полностью открытыми
\end{itemize}
}

% ============================================================
\section{Сравнения}
% ============================================================

% ================================================================
\begin{frame}{Сравнения по модулю $n$}
    \begin{itemize}
        \item Определение: $a \equiv b \pmod n \iff n \mid (a-b)$.
        \item Сравнение по модулю $n$ --- отношение эквивалентности:
        \begin{itemize}
            \item $a\equiv a\pmod n$ (рефлексивность);
            \item $a\equiv b\pmod n \Rightarrow b\equiv a\pmod n$ (симметричность);
            \item $a\equiv b,\ b\equiv c\pmod n \Rightarrow a\equiv c\pmod n$ (транзитивность).
        \end{itemize}
        \item $\mathbb{Z}$ разбивается на $n$ классов эквивалентности (вычетов): $\overline{0}, \overline{1}, \dots, \overline{n-1}$;
        \[\overline{a} = \{a + kn: k\in\mathbb{Z}\}\] 
        \item Сравнения совместимы со сложением и умножением.
        \item Пример: таблицы сложения и умножения по модулю $3$:
        \begin{center}
            \begin{tabular}{c|ccc}
                $+$ & 0 & 1 & 2 \\
                \hline
                0 & 0 & 1 & 2 \\
                1 & 1 & 2 & 0 \\
                2 & 2 & 0 & 1 \\
            \end{tabular}
            \hspace{3em}
            \begin{tabular}{c|ccc}
                $\times$ & 0 & 1 & 2 \\
                \hline
                0 & 0 & 0 & 0 \\
                1 & 0 & 1 & 2 \\
                2 & 0 & 2 & 1 \\
            \end{tabular}
        \end{center}
    \end{itemize}
\end{frame}
\note{%
\begin{itemize}
    \item Совместимость со сложением/умножением --- то, что позволяет вводить арифметику прямо на классах (следующий слайд)
\end{itemize}
}

% ================================================================
\begin{frame}{Мотивация: сравнения как метод исключения решений}
    \begin{itemize}
        \item Если уравнение имеет решение в целых числах, то оно имеет решение и по любому модулю $n$;
        \item Отсутствие решения по какому-то модулю $n$ доказывает отсутствие решения в $\mathbb{Z}$.
        \item Пример: уравнение $x^2+y^2=4k+3$ не имеет целых решений:\\ квадраты по модулю $4$ принимают только значения $0,1$, а их сумма никогда не даёт $3\pmod4$.
        \item Пример: уравнение $x^2-3y^2=2$ не имеет целых решений:\\ квадраты по модулю $3$ равны $0$ или $1$, значение $2$ не достигается.
    \end{itemize}
\end{frame}
\note{%
\begin{itemize}
    \item Это пример "локально-глобального" рассуждения: "локальная" информация (т.е. по модулю $n$) даёт "глобальное" препятствие (т.е. в $\mathbb{Z}$)
    \item Метод дёшев вычислительно и часто быстро закрывает вопрос о разрешимости, когда прямой перебор в $\mathbb{Z}$ бесполезен
    \item Общая идея получит полноценное развитие в курсе позже (тема 1.7 $p$-адические числа)
\end{itemize}
}

% ================================================================
\begin{frame}{Кольцо $\mathbb{Z}/n\mathbb{Z}$}
    \begin{itemize}
        \item Определение: (коммутативное) кольцо с единицей --- множество $R$ с двумя операциями $+,\times$ и выделенными элементами $0,1\in R$, для которых выполнено:
        \begin{itemize}
            \item $(a+b)+c=a+(b+c)$ и $a+b=b+a$ (ассоциативность и коммутативность $+$);
            \item $a+0=a$ для любого $a\in R$ (нейтральный элемент для $+$);
            \item для каждого $a\in R$ существует $-a\in R$ с $a+(-a)=0$ (противоположный для $+$);
            \item $(a\times b)\times c=a\times(b\times c)$ и $a\times b=b\times a$ (ассоциативность и коммутативность $\times$);
            \item $a\times1=a$ для любого $a\in R$ (нейтральный элемент для $\times$);
            \item $a\times(b+c)=a\times b+a\times c$ (дистрибутивность).
        \end{itemize}
        \item $\mathbb{Z}/n\mathbb{Z}$ --- множество классов вычетов с операциями сложения и умножения по модулю $n$ --- коммутативное кольцо с единицей.
        \item $\mathbb{Z}$ --- коммутативное кольцо с единицей.
        \item $\mathbb{Q}$, $\mathbb{R}$ --- кольца, и притом поля: каждый ненулевой элемент обратим для $\times$.
        \item \sagemark
    \end{itemize}
\end{frame}
\note{%
\begin{itemize}
    \item Другие примеры колец и полей рассмотрим в следующей лекции
\end{itemize}
}

% ================================================================
\begin{frame}{Китайская теорема об остатках}
    \begin{itemize}
        \item Если $n = p_1^{e_1}\cdots p_k^{e_k}$, то система сравнений $x\equiv a_i \pmod{p_i^{e_i}}$, $i=1,\dots,k$, имеет единственное решение по модулю $n$.
        \item Пример: $p=3,\ q=5,\ n=15$. Система $x\equiv2\pmod3,\ x\equiv3\pmod5$ имеет единственное решение по модулю $15$: $x=8$;\\ проверка: $8=2\cdot3+2\equiv2\pmod3$, $8=1\cdot5+3\equiv3\pmod5$.
        \item Эквивалентная формулировка --- изоморфизм колец: 
        \[\mathbb{Z}/n\mathbb{Z} \cong \mathbb{Z}/p_1^{e_1}\mathbb{Z}\times\cdots\times\mathbb{Z}/p_k^{e_k}\mathbb{Z},\]
        \item Изоморфизм колец $\psi: R\to S$ --- это биекция, сохраняющая операции: 
        \[\psi(a+b)=\psi(a)+\psi(b),\ \psi(ab)=\psi(a)\psi(b),\ \psi(1_R)=1_S.\]
        \item Частный случай для $n=pq$ (различные простые): 
        \[\mathbb{Z}/n\mathbb{Z}\cong\mathbb{Z}/p\mathbb{Z}\times\mathbb{Z}/q\mathbb{Z}.\]
        \item \sagemark
    \end{itemize}
\end{frame}
\note{%
\begin{itemize}
    \item Доказательство опирается на тождество Безу (существование обратных по модулю каждого $p_i^{e_i}$)
    \item Изоморфизм колец, в частности, сохраняет обратимость элементов: $a$ обратим в $R$ тогда и только тогда, когда $\varphi(a)$ обратим в $S$ --- это пригодится на следующем слайде
\end{itemize}
}

% ================================================================
\begin{frame}{Поле $\mathbb{F}_p$ и функция Эйлера}
    \begin{itemize}
        \item Лемма: элемент $a$ обратим в $\mathbb{Z}/n\mathbb{Z} \Leftrightarrow \gcd(a,n)=1$;
        \item Пример: в $\mathbb{Z}/12\mathbb{Z}$ обратимы $\{1,5,7,11\}$, не обратимы $\{0,2,3,4,6,8,9,10\}$.
        \item Если $n=p$ простое, то $\mathbb{Z}/p\mathbb{Z} = \mathbb{F}_p$ --- поле (обратимость для $\times$).
        \item Малая теорема Ферма: если $p$ простое и $p \nmid a$, то $a^{p-1} \equiv 1 \pmod p$;
        \item Пример: $p=7,\ a=3$: $3^6 = 729 = 104\cdot7+1 \equiv 1\pmod7$.
        \item Функция Эйлера $\varphi(n)$ --- число классов, обратимых по модулю $n$;
        \item Для простого $p$: $\varphi(p)=p-1$;
        \item Для $n=pq$ (различные простые): $\varphi(n)=\varphi(p)\varphi(q)=(p-1)(q-1)$.
        %\item Лемма: $\varphi(n) = n\prod_{p\mid n}\left(1-\tfrac1p\right)$
        \item Теорема Эйлера: если $\gcd(a,n)=1$, то $a^{\varphi(n)} \equiv 1 \pmod n$.
        \item Следствие: если $\gcd(a,n)=1$, то $a^{k\varphi(n)+1} \equiv a \pmod n$ для любого $k\ge0$.
        \item \sagemark
    \end{itemize}
\end{frame}
\note{%
\begin{itemize}
    \item В примере $3\cdot4=12\equiv0$ элементы $3$ и $4$ называются делителями нуля, у делителя нуля обратного быть не может
    \item Критерий обратимости --- прямое следствие тождества Безу: $ax\equiv1\pmod n \iff \exists x,y:\ ax+ny=1 \iff \gcd(a,n)=1$
    \item $\varphi(n)$ --- это в точности порядок мультипликативной группы обратимых элементов $(\mathbb{Z}/n\mathbb{Z})^*$
    \item Функция $\varphi$ --- мультипликативна: $\varphi(mn)=\varphi(m)\varphi(n)$ для взаимно простых $m,n$
    \item Мультипликативная группа обратимых элементов (т.е. ненулевых) $\mathbb{F}_p^*$ всегда циклическая (то есть порождается степенями одного элемента)
    \item \item Следствие потребуется далее для RSA
\end{itemize}
}

% ============================================================
\section{Криптосистема RSA}
% ============================================================

% ================================================================
\begin{frame}{Идея RSA: асимметрия умножения и факторизации}
    \begin{itemize}
        \item Наблюдение: перемножить два больших различных простых $p, q$ легко, а восстановить $p,q$ по $n=pq$ (разложить/факторизовать) --- вычислительно тяжело для больших $n$;
        \item Хотим построить операцию, которую легко сделать в одну сторону (используя знание $p,q$) и трудно обратить, зная только $n$.
    \end{itemize}
\end{frame}
\note{%
\begin{itemize}
    \item Название RSA --- по фамилиям Р. Ривеста, А. Шамира и Л. Адлемана (MIT), опубликовавших конструкцию в 1977--1978 гг. Независимо и раньше (в 1973 г.) эквивалентную схему нашёл К. Кокс в британской GCHQ, но эта работа была засекречена и стала известна только в 1997 г. --- RSA как понятие сформировалось уже вокруг публикации Ривеста--Шамира--Адлемана
    \item Эта асимметрия --- единственная содержательная идея, нужная для RSA (саму криптографическую теорию/атаки не обсуждаем --- только числовая конструкция)
    \item Инструмент для построения такой операции --- как раз теорема Эйлера с предыдущего слайда
\end{itemize}
}

% ================================================================
\begin{frame}{Конструкция RSA}
    \begin{itemize}
        \item Выбираем два различных простых $p, q$, полагаем $n=pq$, $\varphi(n) = (p-1)(q-1)$
        \item Выбираем $e$ взаимно простое с $\varphi(n)$; находим $d$: $ed\equiv1\pmod{\varphi(n)}$ (обратный к $e$ по модулю $\varphi(n)$)
        \item Сообщение $m$: число $0\le m<n$
        \item Шифрование: $m^e \equiv c \pmod n$
        \item Расшифрование: $c^d \equiv m \pmod n$
        \item Корректность: при $\gcd(m,n)=1$: $c^d = m^{ed} = m^{k\varphi(n)+1} \equiv m \pmod n$;
        \item Для $m$ кратных $p$ или $q$ тождество тоже верно (обоснование через КТО и малую теорему Ферма)
        \item \sagemark
    \end{itemize}
\end{frame}
\note{%
\begin{itemize}
    \item если $p=q$, формула $\varphi(n)=(p-1)(q-1)$ неверна: $n=p^2$ и $\varphi(p^2)=p(p-1)$
    \item $d$ существует именно благодаря тождеству Безу (обратимость по модулю $\varphi(n)$)
\end{itemize}
}

% ================================================================
\begin{frame}{RSA на схеме}
    \centering
    \begin{tikzpicture}[
        box/.style={draw=black, thick, rounded corners, align=left, font=\footnotesize, inner sep=8pt},
        collabel/.style={font=\bfseries, align=center},
        arr/.style={-{Latex[length=3mm]}, thick, black},
        darr/.style={-{Latex[length=3mm]}, thick, black, dashed},
        lbl/.style={font=\scriptsize, align=center}
    ]
        \node[collabel, step on=1] (AliceLbl) at (0,5.8) {Алиса};
        \node[collabel, step on=2] (BobLbl) at (8,5.8) {Боб};

        \node[box, step on=1, step draw on=1, text width=5.4cm] (Gen) at (0,4.6) {
            генерирует случайные $p, q$,\\
            считает $n=pq$, $e$, $d$,\\
            ($p, q, d$ --- секрет)
        };
        \node[box, step on=2, step draw on=2, text width=5.4cm] (Bob) at (8,3.9) {
            знает только $n, e$;\\
            шифрует сообщение m:\\
            $c \equiv m^e \bmod n$
        };
        \node[box, step on=3, step draw on=3, text width=5.4cm] (Dec) at (0,2.8) {
            расшифровывает:\\
            $m \equiv c^d \bmod n$\\
            (легко, зная $d, p, q$)
        };
        \node[box, step on=4, step draw on=4, text width=9cm] (Eve) at (4.5,0.5) {
            видит $n, e, c$, но не $p, q, d$\\
            чтобы найти $d$ --- нужно $\varphi(n)$;\\
            вычислить $\varphi(n)$ без $p, q$ $\approx$ факторизовать $n$ (трудно)
        };
        \node[collabel, step on=4, anchor=south] (EveLbl) at ($(Eve.north)+(2,0.2)$) {Ева};

        \def\midx{4} % x координата изгиба -- ровно между колонками Алисы (x=0) и Боба (x=8)

        \coordinate (NEin)  at ($(Bob.west)+(0,0.3)$);  % вход n,e в Bob -- чуть выше середины West-грани
        \coordinate (Cout)  at ($(Bob.west)+(0,-0.3)$); % выход c из Bob -- чуть ниже середины West-грани

        \coordinate (NEbendA) at (\midx,0 |- Gen.east); % изгиб на уровне Gen
        \coordinate (NEbendB) at (\midx,0 |- NEin);     % изгиб на уровне входа в Bob
        \coordinate (CbendA)  at (\midx,0 |- Cout);     % изгиб на уровне выхода из Bob
        \coordinate (CbendB)  at (\midx,0 |- Dec.east); % изгиб на уровне Dec

        \draw[arr, step draw on=2] (Gen.east) -- node[lbl, above, step on=2] {$n, e$} (NEbendA) -- (NEbendB) -- (NEin);
        \draw[arr, step draw on=3] (Cout) -- node[lbl, above, step on=3] {$c$} (CbendA) -- (CbendB) -- (Dec.east);

        \coordinate (NEtap) at ($(NEbendA)+(-0.5,0)$); % точка перехвата n,e -- чуть левее середины
        \coordinate (Ctap)  at ($(CbendA)+(0.5,0)$);   % точка перехвата c -- чуть правее середины

        \draw[darr, step draw on=4] (NEtap) -- (NEtap |- Eve.north);
        \draw[darr, step draw on=4] (Ctap) -- (Ctap |- Eve.north);
    \end{tikzpicture}
\end{frame}
\note{%
\begin{itemize}
    \item Схема Алиса--Боб--Ева:
    \begin{itemize}
        \item Алиса генерирует $p,q$ (легко) и публикует $n=pq,e$;
        \item Боб шифрует сообщение, зная только публичные $n,e$ (легко --- возведение в степень по модулю);
        \item Ева видит $n,e,c$, но не $p,q,d$ --- чтобы прочитать сообщение, ей нужно либо факторизовать $n$, либо иначе найти $d$ (трудно);
        \item Алиса расшифровывает, зная секретные $p,q$ (или сразу $d$) --- снова легко
    \end{itemize}
    \item Это textbook-RSA --- чисто теоретико-числовая конструкция; вопросы стойкости, паддинга, атак --- за рамками нашего курса
\end{itemize}
}

% ================================================================
\begin{frame}{Итоги лекции}
    \begin{itemize}[<.->]
        \item Кольцо $\mathbb{Z}$ и его арифметическая структура;
        \item Сравнения и факторкольцо $\mathbb{Z}/n\mathbb{Z}$;
        \item Первое арифметическое приложение.
        \item В следующий раз --- примеры других колец и их арифметики.
    \end{itemize}
\end{frame}
\note{%
\begin{itemize}
    \item ---
\end{itemize}
}

% ================================================================
\begin{frame}{Литература}
    \begin{itemize}[<.->]
        \item W. Stein, \textit{Elementary Number Theory: Primes, Congruences, and Secrets: A Computational Approach}, 2017 (\href{https://wstein.org/ent/}{wstein.org/ent}). Главы 1--3.
        \item К. Айерлэнд, М. Роузен, Классическое введение в современную теорию чисел, 1987 (Оригинал: K. Ireland, M. Rosen. \textit{A Classical Introduction to Modern Number Theory}, 1990). Главы 1--4.
        \item Н. Коблиц. \textit{Курс теории чисел и криптографии}, 2001 (Оригинал: N. Koblitz, \textit{A Course in Number Theory and Cryptography}, 1994). Глава I.
        \item И.М. Виноградов. \textit{Основы теории чисел}, 1936 (1-е изд., есть современные переиздания). Главы 1---3.
    \end{itemize}
\end{frame}
\note{%
\begin{itemize}
    \item Всё, что было показано сегодня без доказательств (бесконечность простых мы, впрочем, доказали; ОТА, распределение простых, теоремы Ферма/Эйлера, КТО, цикличность $\mathbb{F}_p^*$) --- доказательства и более подробное изложение ищите в источниках выше
    \item Но специально разбираться в доказательствах сейчас не требуется --- курс задуман так, чтобы дать инструменты и интуицию, а полное и строгое изложение теории этих объектов раскрывается в отдельном курсе A-I.
\end{itemize}
}

\end{document}
